% Part: set-theory
% Chapter: ordinals
% Section: introduction

\documentclass[../../../include/open-logic-section]{subfiles}

\begin{document}

\olfileid{sth}{ordinals}{intro}

\olsection{سريزه}

په \olref[z][]{chap} کښې مو د مراتبو د لړۍ د يوه نامتناهي پړاو
شتون د \stagesinf{} په بڼه فرض کړے و (زموږ د نامتناهيت بديهي
اصل هم وګورئ). خو د \stagessucc{} په لرلو سره په نامتناهي پړاو
نۀ شو درېدلے؛ مخته تلل پکار دي. نو د لومړي نامتناهي پړاو نه په
ورپسې پړاو کښې د هغو سټونو ټولې ممکنې ټولګې جوړوو چې په لومړي
نامتناهي پړاو کښې موجودې وې؛ او بيا همدا کار کوو؛ او بيا؛ او بيا؛ \dots

دلته په پټه مو د عدد د يوه «شهودي» تصور کارول پېل کړي دي، چې
له مخې ئې د ټولو طبيعي عددونو \emph{وروسته} هم عددونه کېدلے شي.
په ځانګړې توګه، دلته مفهوم د \emph{نامتناهيت هاخوا ترتيبي عدد}
دے. د دې څپرکي مقصد دا دے چې دا تصور لا کره کړو. د ترتيبي عدد
عمومي مفهوم به وڅېړو، او بيا به په ښکاره ځينې سټونه د خپلو
ترتيبي عددونو په توګه تعريف کړو.

\end{document}
